% SOURCE_STATEMENT: "Can buoyancy reversal remain nonlinearly coupled to external forcings such as radiation?"
% Preparation and final self-audit: CONTRIBUTING.md prompts C and D; source/audit artifacts retained in the local physics_statements_2026_10_11 batch.
% !TeX program = lualatex
\documentclass[11pt]{article}
\usepackage[margin=25mm]{geometry}
\usepackage{fontspec}
\setmainfont{Latin Modern Roman}
\usepackage{amsmath,amssymb,mathtools}
\usepackage{unicode-math}
\setmathfont{Latin Modern Math}
\usepackage{xurl,hyperref}
\hypersetup{colorlinks=true,urlcolor=blue}
\setcounter{secnumdepth}{0}
\setlength{\parindent}{0pt}
\setlength{\parskip}{5pt}
\setlength{\emergencystretch}{3em}
\newcommand{\sref}[2]{\hyperlink{source:#1:#2}{[S#2]}}
\newcommand{\cataloguescope}[1]{\paragraph{Recorded target.}#1}
\newcommand{\R}{\mathbb R}
\newcommand{\C}{\mathbb C}
\newcommand{\N}{\mathbb N}
\newcommand{\Z}{\mathbb Z}
\newcommand{\Q}{\mathbb Q}
\newcommand{\F}{\mathbb F}
\newcommand{\eps}{\varepsilon}
\begin{document}
% UnsolvedMath ID 30001532; source catalogue code OWR-4415-005
\setcounter{section}{30001531}
\section{Nonlinear Coupling of Buoyancy Reversal and Radiation}\label{30001532}
{\small UnsolvedMath ID 30001532 · Mathematical Physics}
\subsection{Definitions and mathematical statement}

\paragraph{Shared notation.}$\mathbb N=\{1,2,\ldots\}$; the standard number fields have their usual meanings. All additional conventions are stated locally.


\paragraph{Definitions and conventions.} Buoyancy reversal is the production of negatively buoyant mixtures during mixing of cloudy and clear air, for example through evaporation near a cloud-top inversion. Buoyancy drives vertical acceleration; radiative heating or cooling supplies an external thermal forcing. Nonlinear coupling means that the combined evolution cannot be accounted for by adding two independent linear responses.

\paragraph{Source scope and qualifications.} The source concerns cloud-top entrainment and specifically lists additional unresolved roles of shear, evaporation rates, droplet settling, and clustering. The extracted coupling question has no single specified PDE or coupling coefficient, so none is invented.

\paragraph{Catalogue mathematical target.}
\paragraph{Catalogue wording retained for traceability.}
Can buoyancy reversal remain nonlinearly coupled to external forcings such as radiation?

\paragraph{Saved literature context.} The supplied catalogue's assessment dated 2026-08-21 records: The extracted question asks whether buoyancy reversal can couple nonlinearly to external forcings such as radiation; it is an ongoing modeling question and the report supplies no definitive answer. The assessment is retained as dated provenance; it is not a new comprehensive literature or open-status certification.
\subsection{Short English statement}
Determine whether evaporative buoyancy reversal couples nonlinearly to radiative and other cloud-top forcings.
\subsection{Sources}
\begin{itemize}
\item[\hypertarget{source:30001532:1}{S1}] ULAM.AI and the UnsolvedMath Contributors, UnsolvedMath, record 30001532 (OWR-4415-005). Catalogue curation under CC BY 4.0; definitions and English exposition added locally.
\url{https://huggingface.co/datasets/ulamai/UnsolvedMath}.
\item[\hypertarget{source:30001532:2}{S2}] Mathematical Theory and Modelling in Atmosphere-Ocean-Science (2010), Oberwolfach Reports, DOI 10.4171/owr/2010/34, pp. 2027--2099; official report source and mathematical context.
\url{https://doi.org/10.4171/owr/2010/34}.
\item[\hypertarget{source:30001532:3}{S3}] Official EMS Press PDF of the same Oberwolfach report; PDF pages 36--37, printed pages 2062--2063.
\url{https://ems.press/content/serial-article-files/46292}.
\end{itemize}
\label{end:30001532}
\end{document}
